\section{Normalization and the immersion reduction}
\label{sec:preliminaries}

All spheres have radius one unless a radius is displayed. We write
\(\sigma_j=\Vol_j(S^j)\) and \(\omega_j=\Vol_j(B^j_1)\), so
\(\sigma_{j-1}=j\omega_j\). For an embedded hypersurface \(\Sigma\),
\(|\Sigma|\) denotes its induced \(m\)-volume. The word closed means
compact without boundary.

The Laplacian is \(\Delta=\operatorname{div}\nabla\). For a two-sided
minimal hypersurface in a sphere, \(x\) is the position vector, \(N\) a
global unit normal, \(A=-dN\) the shape operator, and \(P=|A|^2\).
The Morse index counts the positive eigenvalues, with multiplicity, of
\[
L=\Delta+m+P;
\]
equivalently it is the dimension of a maximal negative subspace for
the area second variation \(-\int_\Sigma uLu\).

\begin{lemma}\label{lem:clifford-ratios}
The numbers \(\vartheta_m=a_m/\sigma_m\) strictly decrease for
integers \(m\ge2\). Moreover,
\[
\vartheta_2=\frac{\pi}{2}<2,
\qquad
\vartheta_4\le\frac32.
\]
Consequently \(a_m<2\sigma_m\) for every \(m\ge2\), and
\(\vartheta_m\le3/2\) for every \(m\ge4\).
\end{lemma}

\begin{proof}
For \(k+l=m\), the product has volume
\[
\sigma_k\sigma_l
\left(\frac{k}{m}\right)^{k/2}
\left(\frac{l}{m}\right)^{l/2}.
\]
Define, for real \(j>0\),
\[
d_j=\sigma_j\left(\frac{j}{2\pi e}\right)^{j/2},
\qquad
\sigma_j=\frac{2\pi^{(j+1)/2}}{\Gamma((j+1)/2)}.
\]
The ratio of the product volume to \(\sigma_m\) is \(d_kd_l/d_{k+l}\).
The usual series for the logarithmic derivatives of \(\Gamma\) gives
\[
\frac{d^2}{dj^2}\log d_j
=\frac1{2j}
-\frac14\sum_{i=0}^{\infty}
\left(i+\frac{j+1}{2}\right)^{-2}>0.
\]
Indeed, strict convexity of \(t\mapsto(t+j/2)^{-2}\) implies
\[
\sum_{i=0}^{\infty}\left(i+\frac{j+1}{2}\right)^{-2}
<
\int_0^\infty(t+j/2)^{-2}\,dt
=\frac2j.
\]
Thus, with \(k>0\) fixed, the derivative in \(l\) of
\(\log(d_kd_l/d_{k+l})\) is negative. Increase the second entry of a
minimizing pair for \(m\) by one to obtain a strictly smaller competitor
for \(\vartheta_{m+1}\). This proves the strict decrease. The stated
values follow from \(k=l=1\) and \(k=l=2\), respectively.
\end{proof}

For a \(d\)-dimensional Euclidean varifold \(V\), its density at \(z\),
when it exists, is
\[
\Theta(V,z)=\lim_{\rho\downarrow0}
\frac{\|V\|(B_\rho(z))}{\omega_d\rho^d}.
\]
Here \(\|V\|\) is its weight measure. We use stationary integral
varifolds with their full integer multiplicities. In particular,
coincident sheets add their masses.

The following multiplicity estimate is due to Nguyen
\cite[Corollary~4.8]{Nguyen2023} and Ge--Li
\cite[Theorem~1.2]{GeLi2022}. We include its cone-monotonicity proof to
make the treatment of source multiplicity explicit.

\begin{proposition}[Immersion reduction]\label{prop:immersion-reduction}
Let \(F:M^m\to S^{m+1}\) be a smooth minimal immersion of a closed
manifold, with \(m\ge2\). If some image point has \(q\) preimages, then
\[
\Vol\bigl(M,F^*g_{\mathrm{round}}\bigr)\ge q\sigma_m.
\]
In particular, every immersion violating the bound in
Theorem~\ref{thm:volume} is an embedding.
\end{proposition}

\begin{proof}
Push forward the parametrized Euclidean cone
\[
(0,\infty)\times M\longrightarrow\mathbb R^{m+2},
\qquad (r,z)\longmapsto rF(z),
\]
with measure \(r^m\,dr\,d\mathrm{vol}_M\). This gives an integral
\((m+1)\)-varifold \(\mathcal C_F\), extended with zero mass at the
origin. Minimality of \(F\) makes it stationary away from the origin.
The boundary term for first variation after truncation at \(r=\epsilon\)
is \(O(\epsilon^m)\); it tends to zero. Thus \(\mathcal C_F\) is
stationary in all of Euclidean space. The corresponding calculation
for general spherical varifolds is given in Lemma~\ref{lem:cone-operation}.

Its centered ball masses are exactly
\[
\|\mathcal C_F\|(B_R(0))
=\frac{R^{m+1}}{m+1}
\Vol\bigl(M,F^*g_{\mathrm{round}}\bigr),
\]
and hence
\[
\Theta(\mathcal C_F,0)
=\frac{\Vol(M,F^*g_{\mathrm{round}})}{\sigma_m}.
\]
Every fiber of a compact immersion is finite: it is compact and its
points are isolated by local embedding neighborhoods. At a nonzero
cone point over an image point with \(q\) preimages, these neighborhoods
give \(q\) smooth local sheets, each of density one. Therefore the density
there is \(q\), even when tangent planes agree or sheets coincide.

For a fixed Euclidean center \(z\), the inclusions
\[
B_{R-|z|}(0)\subset B_R(z)\subset B_{R+|z|}(0)
\]
show that the large-radius mass ratio has limit
\(\Theta(\mathcal C_F,0)\). Stationary monotonicity centered at the
chosen nonzero point consequently gives
\[
q\le \Theta(\mathcal C_F,0).
\]
This proves the inequality. If \(F\) is not injective, take \(q\ge2\)
and use Lemma~\ref{lem:clifford-ratios}. Finally, an injective immersion
from a compact manifold is an embedding.
\end{proof}

\begin{remark}\label{rem:preliminaries-multiplicity}
The preceding argument counts preimages at the actual point; it does
not require transverse self-intersections. On a regular image patch,
the number of coincident sheets is the usual local covering
multiplicity. Distinct minimal sheet germs are analytic and meet in a
set of hypersurface measure zero unless they agree. Thus the generic
multiplicity on such patches also has its ordinary covering
interpretation. In particular, replacing a minimal embedding by a
finite covering multiplies the measured volume by its degree.
\end{remark}

\begin{lemma}\label{lem:fullness}
A closed connected embedded non-equatorial minimal hypersurface
\(\Sigma^m\subset S^{m+1}\) is two-sided and full, meaning that its
position vectors span \(\mathbb R^{m+2}\). It is real analytic, and
\[
\int_\Sigma x=0.
\]
For every nonzero constant vector \(a\), the coordinate function
\(\langle a,x\rangle\) is nontrivial and its zero set has
\(m\)-dimensional measure zero.
\end{lemma}

\begin{proof}
A closed embedded hypersurface in a sphere separates it into two
regions and has a global unit normal. If its image were contained in
a linear hyperplane, it would be an open and closed subset of the
corresponding great sphere, by equality of dimensions. Connectedness
and compactness would make the image that great sphere. This proves
fullness.

The coordinate equation \(\Delta x=-mx\) gives the integral identity.
In analytic ambient coordinates, a minimal hypersurface is locally an
analytic graph by analytic regularity for the minimal graph equation
\cite[Theorem~1.1]{Blatt2024}.
Thus its coordinate functions are analytic. Fullness makes each
\(\langle a,x\rangle\), \(a\ne0\), nontrivial, and a nontrivial analytic
function on a connected manifold has a zero set of measure zero
\cite[Proposition~0]{Mityagin2015}.
\end{proof}
